\section{INTRODUCCIÓN}

La criminalidad urbana y los incidentes que vulneran la seguridad ciudadana constituyen fenómenos sociotécnicos altamente dinámicos, heterogéneos y multidimensionales. Su impacto directo sobre la integridad física de las personas, el bienestar colectivo, la actividad económica y el desarrollo urbano exige que las instituciones de control, monitoreo y planificación formulen estrategias de respuesta y prevención basadas en evidencia empírica rigurosa. En este escenario, la recopilación sistemática de reportes policiales, llamadas de emergencia y crónicas periodísticas ha originado grandes repositorios de datos espaciotemporales no etiquetados.

No obstante, la administración tradicional de la seguridad pública suele operar bajo esquemas reactivos, limitándose al análisis estadístico descriptivo agregado por divisiones político-administrativas rígidas (como distritos, cuadrantes o precintos policiales). Este enfoque tradicional presenta severas limitaciones metodológicas: ignora la continuidad espacial real del territorio, enmascara micro-concentraciones delictivas (\textit{hotspots}) en las fronteras entre zonas, es incapaz de descubrir patrones emergentes de \textit{modus operandi} en descripciones de texto libre y resulta ineficaz para detectar eventos delictivos atípicos o anomalías operativas que se diluyen en el volumen masivo de registros rutinarios.

Frente a este desafío, las técnicas de \textbf{Aprendizaje Automático No Supervisado} (\textit{Unsupervised Machine Learning}) proveen un marco analítico y matemático idóneo para extraer estructuras latentes, segmentaciones intrínsecas y regularidades estadísticas sin requerir supervisión humana ni etiquetas preestablecidas. El presente proyecto desarrolla e implementa un sistema integral de inteligencia urbana y criminología computacional que articula múltiples paradigmas no supervisados:

\begin{itemize}
    \item \textbf{Clustering Espacial Basado en Densidad:} Algoritmos DBSCAN, HDBSCAN y OPTICS formulados con la métrica geodésica de \textit{Haversine} sobre la curvatura de la esfera terrestre, capaces de identificar agrupaciones delictivas de morfología arbitraria y separar formalmente el ruido estocástico de incidentes aislados.
    \item \textbf{Particionamiento Multivariable y Minería Semántica de Texto:} Integración de K-Means con codificación cíclica temporal (seno y coseno de horas y días) y vectorización de texto TF-IDF (\textit{Term Frequency - Inverse Document Frequency}) para clasificar automáticamente narrativas de crímenes y descubrir arquetipos de \textit{modus operandi}.
    \item \textbf{Modelado Probabilístico Continuo:} Mixturas de Gaussianas (GMM) ajustadas mediante el algoritmo Esperanza-Maximización (EM) y elipses de covarianza al 95.4\% ($2\sigma$) para modelar densidades suaves de riesgo.
    \item \textbf{Taxonomía Jerárquica Distrital:} Agrupamiento aglomerativo de enlace Ward que organiza los precintos urbanos según su huella criminal relativa.
    \item \textbf{Reducción de Dimensionalidad y Espacios Latentes:} Proyecciones lineales y no lineales (PCA, TruncatedSVD, PPCA, t-SNE y UMAP) para visualización de alta fidelidad.
    \item \textbf{Detección Neuronal de Anomalías:} Autoencoders profundos con arquitectura de cuello de botella (\textit{bottleneck}) que aíslan crímenes excepcionales mediante el error cuadrático medio de reconstrucción (MSE).
    \item \textbf{Visualización Cartográfica Inmersiva 3D:} Un visor interactivo global WebGL (Globo 3D) que integra de manera dual el conjunto masivo de referencia internacional de Chicago ($+800.000$ registros) con el corpus georreferenciado de seguridad ciudadana de Bolivia ($137$ crónicas estructuradas de los nueve departamentos).
\end{itemize}

El documento expone a continuación el desarrollo conceptual, matemático, metodológico y práctico del sistema, fundamentado bajo el estándar de procesos CRISP-DM / ASUM-DM.

